\section{設計の必然性と依存関係}

本章は規則を導入する前に，どの問題に対してどの構成要素が必要かを分離する．後続章はここで示す
依存方向に従う．

\subsection{MatcherとMatcherSlotを分ける理由}

matcher producerとmatcher consumerを同じ型だけで表すと，coercionを選ぶ根拠が型に現れない．
その場合，特定の構文をcoercion siteとして列挙するか，通常単一化が失敗した後でcoercionを試す
必要がある．前者はfirst-classな関数引数に一般化できず，後者はrollback，solverの探索順，fuelに
受理結果を依存させる．

そこでproducerを$\Matcher{\kappa}{\tau}$，consumer expectationを
$\MatcherSlot{\kappa}{\tau}$とし，checking cutでresolved expected headを観測する．これにより

\begin{center}
\begin{tabular}{p{0.28\linewidth}p{0.29\linewidth}p{0.32\linewidth}}
\toprule
設計上の問題 & 採用する表現 & 得られる性質 \\
\midrule
通常等価性とcoercionの区別
  & expected headの$\mathsf{MatcherSlot}$
  & branch選択がunification失敗に依存しない \\
coercion siteの一般化
  & 構文ではなくchecking judgment
  & $\mathtt{matchAll}$と通常の関数適用が同じ規則を使う \\
producerの保護
  & one-way matcher-to-slot solve
  & consumer demandがproducer capabilityを書き換えない \\
\bottomrule
\end{tabular}
\end{center}

非恒等coercionの終点は常にslotである．matcher-headed expected typeではordinary equalityだけを
許す．expected headが未解決なら構造を推測せず，ordinary branchを選ぶ．

\subsection{三つの推論状態を分ける理由}

demand-directed judgmentは$q;S;\Omega$を左から右へthreadする．三成分は互いに代用できない．

\begin{center}
\begin{tabular}{p{0.12\linewidth}p{0.28\linewidth}p{0.49\linewidth}}
\toprule
状態 & 記録するもの & 独立に必要な理由 \\
\midrule
$q$
  & 二sortの次のfresh番号
  & source判断と実行走査の割当順序を一致させ，変数のboundednessとfuel recursionを証明する \\
$S$
  & そのcutまでのconstraint解
  & 子を左から右へ接続し，expected head，context，一般化を同じchronological stateで解釈する \\
$\Omega$
  & capability variableの生成由来とsolve権限
  & 同じ構文形のvariableでも，rigid，rename-only，structural-flexibleを区別してproducer strengtheningを拒否する \\
\bottomrule
\end{tabular}
\end{center}

$S$だけではcapability variableがscheme producer由来か局所consumer由来かを復元できない．型の
shapeだけから$\Omega$を再計算することもできないため，ledgerはsubstitutionとは別に保持する．一方，
$\Omega$はcoercion demandを決めない．demandはexpected slot head，solve権限はoriginという別軸である．

\subsection{二種類の証明境界}

実行可能推論と関係的source typingは同じpolicyを異なる形で保持する．

\begin{center}
\begin{tabular}{p{0.23\linewidth}p{0.31\linewidth}p{0.35\linewidth}}
\toprule
構成要素 & 必要になる理由 & 後続で可能になること \\
\midrule
raw demand-directed derivation
  & 規則木と三状態の入出力を記録する
  & traversalの構造帰納，state factorization \\
intrinsic Origin certificate
  & 各solveとfreezeがledger policyに従うことを関係的に証明する
  & source側のno-guess／producer protection \\
terminal audit
  & 局所cutでは後続solveに安定しない事実をroot終端で固定する
  & state erasureとvalidator completeness \\
executable trace
  & solver，allocation，event emissionを計算として記録する
  & reconstructionと有限validator \\
terminal validator
  & raw成功のうち全bridge条件を満たすrunだけを公開する
  & 成功等式だけをpremiseとするsoundness \\
$\mathsf{TypingInvariant}$
  & 推論履歴をruntime typingから切り離す
  & preservationとmatching safety \\
\bottomrule
\end{tabular}
\end{center}

Origin certificateは実行traceのコピーではなく，$\mathsf{SourceTyping}$ derivation自身がsolveの
admissibilityを持つために必要である．validatorは逆に，実行runがその関係的条件を満たすことを
有限に検査する．この二つを混同しないため，soundnessはtraceからOriginを再構成し，completenessは
Originとterminal auditから成功traceを再構成する．従ってここでいうcompletenessは，terminal auditを
定義に含む$\mathsf{SourceTyping}$に相対的な主張である．raw derivationが常にauditを満たすことや，
auditと独立に定めた宣言的型付けに対する完全性は含まない．

\subsection{terminal auditを置く三箇所}

局所cutで成立しても，外側の後続substitutionを適用すると再確認が必要な事実がある．すべてのnodeに
任意のfuture suffixへの安定性を要求すると，$\mathtt{let}$のgeneralized setが変わり得るため強すぎる．
そこで公開rootの終端substitutionに固定したauditを，raw derivationと同じ木構造で持つ．

\paragraph{$\mathtt{let}$の局所cutだけを考えた例．}
auditのない$\mathtt{let}$ cutが何を許すかを考える．空のsignature／context，cut substitution
$\mathsf{id}$，value type $\tau_1=\alpha_0\to\alpha_1$に対し，cutでのgeneralizationは
$\sigma_{\mathrm{cut}}=\Gen(\varnothing,\varnothing,\tau_1)=\forall a\,b.\,a\to b$である．bodyの後続solveが
生成したかを問わず，任意の後続suffix $R=[\alpha_1\mapsto\alpha_0]$を許すと，binderは既にcloseされているので
$R\sigma_{\mathrm{cut}}=\sigma_{\mathrm{cut}}$のままである．しかしrootで再計算すべきschemeは
$\Gen(\varnothing,\varnothing,R\tau_1)=\Gen(\varnothing,\varnothing,\alpha_0\to\alpha_0)=\forall a.\,a\to a$となる．従って
$R\sigma_{\mathrm{cut}}\ne\Gen(\varnothing,\varnothing,R\tau_1)$であり，auditなしではdomainとcodomainを独立に
instantiateできる過剰なschemeを残してしまう．$\mathtt{let}$のterminal auditはこの不一致を拒否する．
ただし，これは任意の後続suffixに対して局所cutだけを考えた代数的な例であり，このsuffixを実際のsource traversalが
生成することを示すsource programではない．matcherとpattern constructorでは，同じ再検査原理をterminal
hole evidenceとcapability compatibilityへ適用する．

\begin{center}
\begin{tabular}{p{0.20\linewidth}p{0.34\linewidth}p{0.35\linewidth}}
\toprule
node & 局所cutだけでは足りない事実 & root終端で再確認するもの \\
\midrule
$\mathtt{let}$
  & bodyの後続solveでcontextとvalue typeが変わる
  & 終端context／value typeから再計算したgeneralization \\
matcher literal
  & hole capabilityとfinal capabilityがrenameされ得る
  & evidence，shape，clause capability，exhaustiveness，coverage \\
pattern constructor
  & argument dualとresult capabilityが後続solveを受ける
  & 終端substitution後の$\mathsf{CapCompatible}$ \\
\bottomrule
\end{tabular}
\end{center}

\paragraph{matcher finalizationで実際に到達するterminal auditによる受理損失．}
$\mathsf{List}$用signatureへ，観察不能（opaque）なcapability constructor $\mathsf K$を引数にだけ使う
整形式なdata constructor
\[
\mathtt{consumeK}:\Matcher{\mathsf K}{\mathsf{Int}}\to\mathsf{Witness}
\]
を加える．一要素tupleを分解し，その要素を引数matcherへ委譲する二節matcherを
$\mathtt{fix}\;\mathit{self}\;\mathit{argument}$で包むと，次の閉じたlambdaが得られる．
\[
\lambda m.\;\bigl(
  \mathtt{consumeK}\;m,
  (\mathtt{fix}\;\mathit{self}\;\mathit{argument}.\,
    \mathtt{matcher}[\mathtt{tuple}_1\mapsto\mathit{argument};\mathtt{catch}])\;m
\bigr).
\]
$\mathit{self}$は使わず，$\mathtt{fix}$は引数slotの局所capability variableを作るためだけに用いる．
整形式signatureの下で$\mathsf{inferRaw}$は成功し，raw targetは
\[
\Matcher{\mathsf K}{\mathsf{Int}}\to
\bigl(\mathsf{Witness}\times
  \Matcher{\prodty(\mathsf K)}{\prodty(\mathsf{Int})}\bigr)
\]
である．9個のterminal checkは順に
\[
(\mathtt{true},\mathtt{true},\mathtt{true},\mathtt{true},\mathtt{true},
 \mathtt{true},\mathtt{true},\mathtt{false},\mathtt{true})
\]
となり，8番目の$\mathsf{traceFinalizationSuffixCheck}$だけが失敗して公開$\mathsf{infer}$は拒否する．
matcher finalizationは，引数slotから借りたcapability variableをfreeze対象から除く．局所検査後の
applicationがこのvariableを$\mathsf K$へ特殊化すると，終端検査は委譲先の$\mathsf K$まで再帰的に
観察しようとして拒否する．しかしmatcher自身が観察するのは外側の一要素tupleだけであり，要素の観察は
引数matcherへ委譲される．これはrawでは受理されるが公開推論では拒否されるという，terminal auditによる
具体的な受理損失である．

形式化では，compile-time guardがraw成功，公開拒否，型，9検査の結果を計算し，通常の定理が任意の
raw成功等式から対応するdemand-directed derivationを構成する．さらに閉lambdaがclosureへ評価されることを
証明している．auditと独立な宣言的型付けと，それに基づく安全性定理はまだないため，
このprogramを一般的な意味で安全とする定理や，auditによる損失の正確な範囲は未確立である．修正時には，
matcher自身が観察した構造と，引数slotから完全な値として借りたcapabilityの由来をfinalizationまで保つ必要がある．
借りたvariableも一律にfreezeすると意図した後続特殊化を失い，すべてのopaque capabilityを許すとmatcher自身が
opaqueな内部構造を観察する場合まで受理し得る．

variable nodeは追加auditを持たない．canonical schemeのfinite openingとsubstitution transportから
終端lookupとの一致を代数的に導く．

\subsection{公開定理までの依存グラフ}

証明は局所代数から公開定理へ次の順で積み上げる．

\begin{center}
\begin{tabular}{c}
exact solver + ledger admissibility + canonical scheme laws \\
$\Downarrow$ \\
全raw familyのstate extension + chronological factorization + solved-form保存 \\
$\Downarrow$ \\
Origin tree + terminal audit + executable event coverage \\
$\Downarrow$ \\
\begin{tabular}{c@{\qquad}c@{\qquad}c}
infer $\Rightarrow$ SourceTyping & SourceTyping $\Rightarrow$ infer & SourceTyping $\Rightarrow$ TypingInvariant
\end{tabular} \\
$\Downarrow$ \\
受理同値・決定可能性・target一意性
\qquad
SourceTypingによる公開安全性
\end{tabular}
\end{center}

$\mathsf{FrozenSigWF}$はsource-facingな唯一のglobal signature条件である．受理完全性にはterminal
factsを実行stateへ輸送するclosednessとcanonical arm checkerを与え，安全性にはscheme closednessと
動的整合性を与える．低レベル補題は必要なfieldを明示してよいが，公開callerへ別premiseとして
露出しない．
